% Figure: the differential control volume from which the heat equation is
% derived. A slab element of thickness dx carries a conducted flux q(x) in at
% the left face and q(x+dx) out at the right, with internal generation and a
% stored-energy term accounting for the imbalance. The Taylor expansion of the
% flux is the source of the second-derivative term in the heat equation.
\begin{figure}[htbp]
\centering
\begin{tikzpicture}[every node/.style={font=\small}]
% the control volume
\draw[fill=engblue!8, draw=engblue, thick] (0,0) rectangle (3.2,2.4);
% left and right face markers
\draw[engred, thick] (0,-0.05) -- (0,2.45);
\draw[engred, thick] (3.2,-0.05) -- (3.2,2.45);
% incoming flux arrow
\draw[engred, -{Stealth}, very thick] (-1.8,1.2) -- (-0.05,1.2);
\node[engred, font=\scriptsize, anchor=south] at (-0.9,1.25) {$\dot{Q}_x$};
% outgoing flux arrow
\draw[engred, -{Stealth}, very thick] (3.25,1.2) -- (5.0,1.2);
\node[engred, font=\scriptsize, anchor=south] at (4.1,1.25) {$\dot{Q}_{x+dx}$};
% generation
\node[engorange, font=\scriptsize] at (1.6,1.7) {$\dot{q}_v\,A\,dx$};
\draw[engorange, -{Stealth}] (1.6,1.5) -- (1.6,1.05);
% storage
\node[enggray, font=\scriptsize] at (1.6,0.6) {$\rho c_p A\,dx\,\dfrac{\partial T}{\partial t}$};
% dx dimension
\draw[<->, enggray] (0,-0.45) -- (3.2,-0.45);
\node[enggray, font=\scriptsize, anchor=north] at (1.6,-0.5) {$dx$};
% area label
\node[engblue, font=\scriptsize, anchor=south] at (1.6,2.5) {cross-section $A$};
% x axis
\draw[-{Stealth}] (-1.8,-0.9) -- (5.0,-0.9);
\node[anchor=north] at (5.0,-0.95) {$x$};
\end{tikzpicture}
\caption[The differential control volume for the heat equation]{The
differential control volume from which the heat equation follows. Heat is
conducted in at the left face at rate $\dot{Q}_x$ and out at the right face at
rate $\dot{Q}_{x+dx}$, generation adds $\dot{q}_v A\,dx$ inside, and the
imbalance is stored as the rate of change of internal energy $\rho c_p A\,dx\,
\partial T/\partial t$. Expanding the outgoing flux as $\dot{Q}_{x+dx} =
\dot{Q}_x + (\partial \dot{Q}/\partial x)\,dx$ and substituting Fourier's law
$\dot{Q} = -kA\,\partial T/\partial x$ turns the balance into the
one-dimensional heat equation. The second spatial derivative enters because the
stored energy is set by the net flux difference across the element, not by the
flux itself.}
\label{fig:vol04ch05:control-volume}
\end{figure}
